\documentclass[a4paper,11pt]{letter}
\usepackage[default]{sourcesanspro}
\usepackage{geometry}
\usepackage{setspace}
\usepackage[hidelinks]{hyperref}

\geometry{a4paper, top=3cm, bottom=3cm, left=2.5cm, right=2.5cm}

\address{
    \textbf{Yuhang Chen}\\
    Independent Researcher\\
    Zhejiang, China\\
    \href{mailto:yuhang.chen200203@gmail.com}{yuhang.chen200203@gmail.com}
}

\begin{document}

\begin{letter}{The Editor\\
\textit{Nonlinear Dynamics}}

\opening{Dear Editor,}

I am submitting my manuscript entitled \textbf{``Emergence of Stable Cohesive Structures in Sparse Dynamic Networks''} for consideration in \textit{Nonlinear Dynamics}.

This paper makes a novel theoretical contribution: we define and rigorously prove the emergence of a new graph-theoretic structure---the $\delta$-cohesive partition---in Hopfield-Hebbian dynamical systems on sparse networks. The work directly fills a gap left by Centorrino, Bullo and Russo (Automatica, 2024), who proved global stability of such systems but did not characterize the structure of the converged state.

The core results are three theorems forming a deductive chain. Theorem~1 constructs an exact Lyapunov function whose derivative decomposes into two independent sums of non-positive squares. Theorem~2 uses LaSalle's invariance principle to prove global convergence to equilibrium for any finite undirected graph. Theorem~3 proves that every equilibrium exhibits a $\delta$-cohesive partition with a sharp, analytically computable weight gap $\delta = 2/\alpha$.

Beyond the core theorems, the paper provides:
\begin{itemize}
    \item A \textbf{bifurcation analysis} of the first-order phase transition at a critical learning rate, identifying the saddle-node mechanism responsible for the emergence of cohesive structure;
    \item Three \textbf{graph-theoretic applications}: $k$-plex dense subgraph formation, 100\% structural balance on the induced signed graph, and spectral correspondence with the graph Laplacian Fiedler vector;
    \item A \textbf{1000-node three-dimensional spatial network} simulation confirming dimension-independence of the emergence mechanism;
    \item \textbf{Full Lean~4 formal verification} of all three theorems (10 modules, 3,302 proof jobs, zero \texttt{sorry}).
\end{itemize}

I believe this manuscript aligns with \textit{Nonlinear Dynamics}'s scope in nonlinear dynamical systems, bifurcation theory, and complex network dynamics. The work provides the first rigorous Lyapunov-based bridge between stability theory and graph-theoretic self-organization.

I am an independent researcher. This is my first submission to a peer-reviewed journal in mathematics. I confirm that this manuscript is original, has not been published previously, and is not currently under consideration for publication elsewhere. I declare no conflicts of interest.

Thank you for your time and consideration.

\vspace{2em}
\noindent Sincerely,\\[3em]
\textbf{Yuhang Chen}\\
Independent Researcher

\end{letter}
\end{document}
